% ============================================================
% CHAPITRE 5 : Stabilité, Marges, Routh, Nyquist
% ============================================================
\chapterbanner{Chapitre 5 --- Stabilit\'e \& Marges}{ch5red}

\begin{tcolorbox}[colback=ch5red!4,colframe=ch5red!80!black,title={Stabilit\'e BIBO}]
\tightmath{
\text{Stable (BIBO)} \Leftrightarrow \textbf{tous les p\^oles} : \text{Re}(s_i) < 0
}
\begin{itemize}
  \item Marginalement stable : p\^ole(s) simple(s) sur axe $j\omega$
  \item Instable : au moins un p\^ole avec Re$(s_i) > 0$
\end{itemize}
\textbf{Condition n\'ecessaire (Hurwitz)} : tous les $a_i$ du polyn\^ome caract\'eristique de \textbf{m\^eme signe} et \textbf{non nuls}.
\piege{N\'ecessaire mais PAS suffisante pour $n \geq 3$. Pour $n \leq 2$, elle est aussi suffisante.}
\end{tcolorbox}

\begin{tcolorbox}[colback=ch5red!2,colframe=ch5red!80!black,title={Crit\`ere de Routh-Hurwitz}]
Pour $D(s) = a_n s^n + a_{n-1}s^{n-1} + \cdots + a_0$ :
\tightmath{
\begin{array}{c|cccc}
s^n & a_n & a_{n-2} & a_{n-4} & \cdots \\
s^{n-1} & a_{n-1} & a_{n-3} & a_{n-5} & \cdots \\
s^{n-2} & b_1 & b_2 & b_3 & \cdots \\
s^{n-3} & c_1 & c_2 & \cdots &
\end{array}
}
\tightmath{
b_1 = \frac{a_{n-1} a_{n-2} - a_n a_{n-3}}{a_{n-1}} \quad b_2 = \frac{a_{n-1} a_{n-4} - a_n a_{n-5}}{a_{n-1}}
}
\tightmath{
c_1 = \frac{b_1 a_{n-3} - a_{n-1} b_2}{b_1}
}
\tightmath{
\boxed{\text{Nb p\^oles instables} = \text{nb changements de signe, 1\textsuperscript{\`ere} colonne}}
}
\vars{Stable $\Leftrightarrow$ tous les \'el\'ements 1\textsuperscript{\`ere} col. $> 0$. Z\'ero en 1\textsuperscript{\`ere} col. $\to$ remplacer par $\varepsilon > 0$. Ligne nulle $\to$ polyn\^ome auxiliaire $dP/ds$.}
\end{tcolorbox}

\begin{tcolorbox}[colback=ch5red!4,colframe=ch5red!80!black,title={Exemples Routh rapides}]
\textbf{Ordre 2} : $s^2 + a_1 s + a_0 = 0$. Stable $\Leftrightarrow a_1 > 0$ et $a_0 > 0$.

\textbf{Ordre 3} : $s^3 + a_2 s^2 + a_1 s + a_0 = 0$ :
\tightmath{
\boxed{\text{Stable} \Leftrightarrow a_2 > 0,\ a_0 > 0,\ a_2 a_1 > a_0}
}

\textbf{Ordre 4} : $s^4 + a_3 s^3 + a_2 s^2 + a_1 s + a_0 = 0$ :\\
$a_3 > 0$, $a_0 > 0$, $b_1 = a_3 a_2 - a_1 > 0$, et $\boxed{a_1(a_3 a_2 - a_1) > a_3^{2}\, a_0}$ (i.e. $b_1 a_1 > a_3^{2} a_0$).
\end{tcolorbox}

\begin{tcolorbox}[colback=ch5red!2,colframe=ch5red!80!black,title={Stabilit\'e en boucle ferm\'ee}]
\'Equation caract\'eristique :
\tightmath{
\boxed{1 + G_o(s) = 0} \quad G_o = G_c\,G_a
}
BF stable $\Leftrightarrow$ toutes les racines de $1 + G_o(s) = 0$ ont Re $< 0$.
\end{tcolorbox}

\begin{tcolorbox}[colback=ch5red!4,colframe=ch5red!80!black,title={Marge de gain ($A_m$)}]
$\omega_\pi$ = \textbf{croisement de phase} : $\angle G_o(\jw_\pi) = -180°$.
\tightmath{
\boxed{A_m = \frac{1}{|G_o(\jw_\pi)|}} \quad \boxed{A_{m,\text{dB}} = -|G_o(\jw_\pi)|_{\text{dB}}}
}
\vars{Sur Bode : $A_{m,\text{dB}}$ = distance module $\leftrightarrow$ 0~dB, \textbf{lue \`a $\omega_\pi$}. Stable si courbe \textbf{en dessous} de 0~dB. Typ. : $A_m \geq 6$--$10$ dB.}
\end{tcolorbox}

\begin{tcolorbox}[colback=ch5red!2,colframe=ch5red!80!black,title={Marge de phase ($\varphi_m$)}]
$\omega_{co}$ = \textbf{croisement de gain} : $|G_o(\jw_{co})| = 1$ (0~dB).
\tightmath{
\boxed{\varphi_m = 180° + \angle G_o(\jw_{co})}
}
\vars{Sur Bode : $\varphi_m$ = distance phase $\leftrightarrow$ $-180°$, \textbf{lue \`a $\omega_{co}$}. Stable si phase \textbf{au-dessus} de $-180°$. Objectif : $30° < \varphi_m < 60°$.}
\end{tcolorbox}

\begin{tcolorbox}[colback=ch5red!4,colframe=ch5red!80!black,title={Lecture sur le diagramme de Bode},top=0.4mm,bottom=0.4mm,before skip=1pt,after skip=1pt]
\begin{center}
\begin{tikzpicture}[scale=0.56, transform shape, >=Stealth]
  \begin{scope}
    \draw[->] (0,0) -- (5.5,0) node[right]{$\log\omega$};
    \draw[->] (0,-2) -- (0,2.2) node[above]{dB};
    \draw[dashed, gray] (0,0) -- (5.5,0); \node[left] at (0,0) {0};
    \draw[thick, ch5red!80!black] (0,1.8) -- (1.5,1.2) -- (2.5,0) -- (4,-1.5) -- (5.5,-2.5);
    \draw[red, thick, dashed] (2.5,-2) -- (2.5,0.3); \node[red, above] at (2.5,0.4) {$\omega_{co}$};
    \draw[orange, thick, dashed] (4,-2) -- (4,0.2); \node[orange, above] at (4,0.3) {$\omega_\pi$};
    \draw[<->, green!60!black] (4,-1.5) -- (4,0); \node[right, green!60!black] at (4.1,-0.7) {$A_m$};
    \node[align=left] at (1.35,2) {$\alpha$ : pente BF\\[-1pt]$/(-20)$};
    \node[align=left] at (4.25,-2.62) {$d$ : pente HF\\[-1pt]$/(-20)$};
  \end{scope}
  \begin{scope}[xshift=6.8cm]
    \draw[->] (0,0) -- (5.5,0) node[right]{$\log\omega$};
    \draw[->] (0,-2.2) -- (0,0.5) node[above]{$\angle$};
    \draw[dashed, gray] (0,-1.57) -- (5.5,-1.57); \node[left] at (0,-1.57) {$-180°$};
    \draw[thick, ch5red!80!black] (0,-0.3) -- (1.5,-0.7) -- (2.5,-1.1) -- (4,-1.57) -- (5.5,-2.1);
    \draw[red, thick, dashed] (2.5,-2.2) -- (2.5,-0.8); \node[red, above] at (2.5,-0.75) {$\omega_{co}$};
    \draw[orange, thick, dashed] (4,-2.2) -- (4,-1.35); \node[orange, above] at (4,-1.3) {$\omega_\pi$};
    \draw[<->, green!60!black] (2.5,-1.57) -- (2.5,-1.1); \node[left, green!60!black] at (2.3,-1.3) {$\varphi_m$};
  \end{scope}
\end{tikzpicture}
\end{center}
\vars{$\alpha$ : pente BF du module ; $d$ : pente HF. $\omega_{co} \Rightarrow \varphi_m$, $\omega_\pi \Rightarrow A_m$.}
\end{tcolorbox}

\begin{tcolorbox}[colback=ch5red!2,colframe=ch5red!80!black,title={Relations approx. $\varphi_m$ -- $\zeta$ -- $D_\%$}]
\tightmath{
\boxed{\zeta \approx \frac{\varphi_m}{100°}} \quad (\text{pour } 20° < \varphi_m < 70°)
}
\vars{$\varphi_m = 45° \Rightarrow \zeta \approx 0.45 \Rightarrow D_\% \approx 20\%$. $\varphi_m = 65° \Rightarrow \zeta \approx 0.65 \Rightarrow D_\% \approx 7\%$.}
\end{tcolorbox}

\begin{tcolorbox}[colback=ch5red!4,colframe=ch5red!80!black,title={Approximations BF depuis Bode}]
\textbf{$T_{\text{r\'eg}}$ en BF} (approx. demi-p\'eriode \`a $\omega_{co}$) :
\tightmath{
\boxed{T_{\text{r\'eg}} \approx \frac{\pi}{\omega_{co}}}
}
\textbf{Retard pur marginalisant la BF} (consume $\varphi_m$ : $\omega_{co} T_r = \varphi_{m,\text{rad}}$) :
\tightmath{
\boxed{T_r^{\max} = \frac{\varphi_{m,\text{rad}}}{\omega_{co}} = \frac{\varphi_m° \cdot \pi}{180° \cdot \omega_{co}}}
}
\piege{$\varphi_m$ en \textbf{radians} : $\varphi_{m,\text{rad}} = \varphi_m° / 57.3$. Ex. : $\varphi_m{=}40°$, $\omega_{co}{=}600$ $\Rightarrow$ $T_r^{\max} = \frac{40\pi}{180\cdot600} \approx 1.2\,\text{ms}$.}
\end{tcolorbox}

\begin{tcolorbox}[colback=ch5red!4,colframe=ch5red!80!black,title={Crit\`ere de Nyquist}]
\textbf{Cas simplifi\'e (BO stable, $P{=}0$)} :
BF stable $\Leftrightarrow$ Nyquist $G_o(\jw)$ \textbf{n'entoure pas} $(-1, 0)$.

\textbf{Cas g\'en\'eral :}
\tightmath{
\boxed{Z = N + P}
}
\begin{itemize}
  \item $Z$ = nb p\^oles BF instables (on veut $Z{=}0$)
  \item $N$ = nb encerclements \textbf{horaires} de $(-1,0)$
  \item $P$ = nb p\^oles BO instables (connu)
\end{itemize}
Stable $\Leftrightarrow$ $N = -P$ (encerclements \textbf{anti-horaires}).
\piege{BO stable ($P{=}0$) : $N{=}0$. BO instable ($P{>}0$) : il faut $P$ encerclements anti-horaires.}
\end{tcolorbox}

\begin{tcolorbox}[colback=ch5red!2,colframe=ch5red!80!black,title={Lieu des racines (rappel)}]
Lieu des p\^oles BF quand $K$ varie ($0 \to \infty$) :
\begin{itemize}
  \item D\'epart : p\^oles de BO. Arriv\'ee : z\'eros de BO ou $\infty$
  \item Nb branches $= n$. Asymptotes : angle $= \frac{(2k{+}1)\cdot 180°}{n-m}$
  \item Centre : $\sigma_a = \frac{\sum p_i - \sum z_i}{n - m}$
  \item Sur axe r\'eel : \`a droite d'un nombre \textbf{impair} de p\^oles+z\'eros
  \item Instabilit\'e quand une branche traverse l'axe $j\omega$
\end{itemize}
\end{tcolorbox}

\begin{tcolorbox}[colback=ch5red!4,colframe=ch5red!80!black,title={Proc\'edure stabilit\'e param\'etrique en 5 \`etapes}]
\begin{enumerate}
  \item Former l'\'equation caract\'eristique : $1 + K\,G_o(s) = 0$
  \item Mettre sous forme polynomiale : $a_n(K)s^n + \cdots + a_0(K)=0$
  \item Construire Routh avec coefficients d\'ependant de $K$
  \item Imposer 1\textsuperscript{\`ere} colonne $>0$ et r\'esoudre les in\'egalit\'es
  \item Croiser avec les marges Bode/Nyquist pour garder une zone robuste
\end{enumerate}
\tightmath{
\text{Zone utile en pratique : } \varphi_m\in[35^\circ,60^\circ],\quad A_m\ge 6\text{ dB}
}
\astuce{Si plusieurs intervalles de $K$ sont math\'ematiquement stables, choisir celui qui respecte aussi la rapidit\'e et le d\'epassement demand\'es.}
\end{tcolorbox}
